% LaTeX Template for AI Problem Analysis Output
% AMC 12A 2017 Problem 12: Strategic Analysis
% Structure your analysis results using this template with colorful boxes

\documentclass[]{article}
\usepackage[utf8]{inputenc}
\usepackage{amsmath, amssymb, amsthm}
\usepackage{geometry}
\usepackage{xcolor}
\usepackage[most]{tcolorbox}
\usepackage{enumitem}
\usepackage{fancyhdr}

% Page setup
\geometry{letterpaper, margin=1in}
\pagestyle{fancy}
\fancyhf{}
\rhead{AMC 12A 2017 - Problem 12}
\lhead{Strategic Analysis}
\rfoot{Page \thepage}

% Custom colored boxes
\newtcolorbox{problemconfigbox}[1][]{
    colback=blue!5!white,
    colframe=blue!75!black,
    fonttitle=\bfseries,
    title={1. PROBLEM CONFIGURATION ANALYSIS},
    #1
}

\newtcolorbox{strategicbox}[1][]{
    colback=red!5!white,
    colframe=red!75!black,
    fonttitle=\bfseries,
    title={2. STRATEGIC FRAMEWORK APPLICATION},
    #1
}

\newtcolorbox{tacticalbox}[1][]{
    colback=green!5!white,
    colframe=green!75!black,
    fonttitle=\bfseries,
    title={3. TACTICAL METHODOLOGY SELECTION},
    #1
}

\newtcolorbox{conversionbox}[1][]{
    colback=yellow!5!white,
    colframe=yellow!75!black,
    fonttitle=\bfseries,
    title={4. CONVERSION TECHNIQUES APPLICATION},
    #1
}

\newtcolorbox{verificationbox}[1][]{
    colback=orange!5!white,
    colframe=orange!75!black,
    fonttitle=\bfseries,
    title={5. VERIFICATION STRATEGY DESIGN},
    #1
}

\newtcolorbox{roadmapbox}[1][]{
    colback=cyan!5!white,
    colframe=cyan!75!black,
    fonttitle=\bfseries,
    title={6. SOLUTION ROADMAP},
    #1
}

\newtcolorbox{competitionbox}[1][]{
    colback=purple!5!white,
    colframe=purple!75!black,
    fonttitle=\bfseries,
    title={7. COMPETITION STRATEGY INTEGRATION},
    #1
}

\newtcolorbox{keyinsightbox}[1][]{
    colback=teal!10!white,
    colframe=teal!75!black,
    fonttitle=\bfseries,
    title={KEY INSIGHT},
    #1
}

\newtcolorbox{warningbox}[1][]{
    colback=red!10!white,
    colframe=red!75!black,
    fonttitle=\bfseries,
    title={WARNING - Common Pitfall},
    #1
}

\newtcolorbox{tipbox}[1][]{
    colback=green!10!white,
    colframe=green!75!black,
    fonttitle=\bfseries,
    title={PRO TIP},
    #1
}

\title{\textbf{Strategic Analysis: AMC 12A 2017 Problem 12}\\
\large{Horses on a Circular Track - Divisibility and LCM Problem}}
\author{Comprehensive Framework Analysis}
\date{\today}

\begin{document}

\maketitle

\section*{Problem Statement}
There are $10$ horses, named Horse 1, Horse 2, $\ldots$, Horse 10. They get their names from how many minutes it takes them to run one lap around a circular race track: Horse $k$ runs one lap in exactly $k$ minutes. At time 0 all the horses are together at the starting point on the track. The horses start running in the same direction, and they keep running around the circular track at their constant speeds. The least time $S > 0$, in minutes, at which all $10$ horses will again simultaneously be at the starting point is $S = 2520$. Let $T>0$ be the least time, in minutes, such that at least $5$ of the horses are again at the starting point. What is the sum of the digits of $T$?

\textbf{Answer Choices:}
$\textbf{(A)}\ 2\qquad\textbf{(B)}\ 3\qquad\textbf{(C)}\ 4\qquad\textbf{(D)}\ 5\qquad\textbf{(E)}\ 6$

\vspace{1em}

\begin{problemconfigbox}
\subsection*{A. Problem Classification}
\begin{itemize}[leftmargin=*]
    \item \textbf{Problem Type}: To Find - determine the minimum time $T$ when at least 5 horses return to start
    \item \textbf{Difficulty Band}: Mid-band (Problem 12 of 25) - requires divisibility/LCM insight but not advanced techniques
    \item \textbf{Competition Context}: AMC12 (75 min, 25 problems) - approximately 3 minutes allocated per problem
    \item \textbf{Mathematical Domain}: Number Theory (divisibility, LCM) with combinatorial aspects (counting divisors)
\end{itemize}

\subsection*{B. Problem Structure Identification}
\begin{itemize}[leftmargin=*]
    \item \textbf{Given Information}: 
    \begin{itemize}
        \item 10 horses with lap times: Horse $k$ completes a lap in $k$ minutes ($k = 1, 2, \ldots, 10$)
        \item All horses start at time 0 from the same starting point
        \item $S = 2520$ is the least time when all 10 horses are simultaneously at start (verification: $\text{lcm}(1,2,\ldots,10) = 2520$)
    \end{itemize}
    \item \textbf{Constraints}: 
    \begin{itemize}
        \item $T > 0$ (time must be positive)
        \item At least 5 horses must be at starting point
        \item Seeking the \textit{minimum} such time
    \end{itemize}
    \item \textbf{Objective}: Find the minimum time $T$ when at least 5 horses are at the starting point, then compute the sum of digits of $T$
    \item \textbf{Hidden Assumptions}: 
    \begin{itemize}
        \item Horse $k$ is at the starting point at time $t$ if and only if $k$ divides $t$
        \item We need the smallest positive integer $t$ that has at least 5 divisors from the set $\{1, 2, 3, \ldots, 10\}$
    \end{itemize}
    \item \textbf{Answer Choice Leverage}: Small answer choices (2-6) suggest $T$ is a small number with easily computable digit sum. The problem essentially reduces to finding a number with specific divisibility properties.
\end{itemize}

\subsection*{C. Strategic Orientation}
\begin{itemize}[leftmargin=*]
    \item \textbf{Hypothesis}: At time $t$, Horse $k$ is at the starting point $\Leftrightarrow$ $k \mid t$
    \item \textbf{Conclusion}: Find the smallest positive integer $T$ such that $|\{k \in \{1,2,\ldots,10\} : k \mid T\}| \geq 5$
    \item \textbf{Notation}: Let $d(n)$ denote the number of divisors of $n$ from the set $\{1,2,\ldots,10\}$
    \item \textbf{Intuitive Approach}: 
    \begin{itemize}
        \item Check small positive integers systematically
        \item Look for numbers with many small divisors
        \item Products of small primes (like $2^a \cdot 3^b$) tend to have many divisors
    \end{itemize}
    \item \textbf{Cross-Domain Potential}: 
    \begin{itemize}
        \item Pure Number Theory: Divisor counting
        \item Combinatorics: Counting divisors from a specific set
        \item Computational: Systematic enumeration for small numbers
    \end{itemize}
\end{itemize}
\end{problemconfigbox}

\begin{strategicbox}
\subsection*{A. Psychological Strategies}
\begin{itemize}[leftmargin=*]
    \item \textbf{Mental Toughness}: This is a mid-band problem where the key insight (divisibility condition) must be identified quickly. Don't overthink - trust the systematic approach
    \item \textbf{Creativity Cultivation}: Recognize that the problem is asking for divisor properties rather than explicit simulation of horse positions
    \item \textbf{Iterative Refinement}: Start with small candidates and systematically increase until finding the first valid answer
\end{itemize}

\subsection*{B. Investigation Strategies}
\begin{itemize}[leftmargin=*]
    \item \textbf{Startup Orientation}: Immediately recognize this is a divisibility problem - Horse $k$ returns to start when $k \mid t$
    \item \textbf{Get Your Hands Dirty}: Test small numbers: $t = 1, 2, 3, \ldots$ and count divisors
    \item \textbf{Penultimate Step}: Once we find $T$ with at least 5 divisors, compute digit sum (the actual answer)
    \item \textbf{Wishful Thinking}: "What if $T$ were very small, like 12 or 16?" Then check these candidates first
    \item \textbf{Make It Easier}: Instead of tracking 10 horses, just count how many numbers from $\{1,2,\ldots,10\}$ divide various small integers
    \item \textbf{Conjecture via Small Cases}: 
    \begin{itemize}
        \item $t=1$: divisors $\{1\}$ - only 1 horse
        \item $t=2$: divisors $\{1,2\}$ - 2 horses
        \item $t=6$: divisors $\{1,2,3,6\}$ - 4 horses
        \item $t=12$: divisors $\{1,2,3,4,6\}$ - 5 horses! (Check: $12 = 2^2 \cdot 3$)
    \end{itemize}
    \item \textbf{Visualization}: Number line showing which horses are at start at different times
\end{itemize}

\subsection*{C. Late-Band Playbook (If Applicable)}
\begin{itemize}[leftmargin=*]
    \item \textbf{Applicability}: Problem 12 is mid-band, not late-band, so these techniques are secondary
    \item \textbf{Pattern Recognition}: Numbers with many divisors typically have the form $2^a \cdot 3^b$ or $p^k$ for small primes
    \item \textbf{Answer Choice Working Backwards}: The digit sums 2-6 correspond to small numbers (e.g., digit sum 3 could be 12, 21, 30, 111, etc.)
\end{itemize}

\subsection*{D. Argument Construction}
\begin{itemize}[leftmargin=*]
    \item \textbf{Direct Construction}: Systematically check $t = 1, 2, 3, \ldots$ until finding minimum with 5+ divisors from $\{1,\ldots,10\}$
    \item \textbf{Existence Proof}: We know at least one solution exists ($t = 2520$ has all 10 divisors), so minimum definitely exists
\end{itemize}
\end{strategicbox}

\begin{tacticalbox}
\subsection*{A. Core Mathematical Tactics}
\begin{itemize}[leftmargin=*]
    \item \textbf{Primary Tactic}: \textit{Divisibility and counting divisors}
    \begin{itemize}
        \item Horse $k$ returns to start at time $t$ $\Leftrightarrow$ $k \mid t$
        \item Count divisors of $t$ from the set $\{1,2,\ldots,10\}$
    \end{itemize}
    \item \textbf{Supporting Tactic}: \textit{Systematic enumeration}
    \begin{itemize}
        \item Check small numbers in order
        \item Numbers with form $2^a \cdot 3^b$ are good candidates (many small divisors)
    \end{itemize}
    \item \textbf{Extreme Principle}: Start from the smallest possible time ($t=1$) and work upward
    \item \textbf{Invariant/Monovariant}: As $t$ increases, the number of divisors doesn't necessarily increase monotonically, but we're searching for the minimum $t$ satisfying our condition
\end{itemize}

\subsection*{B. Specialized Techniques}
\begin{itemize}[leftmargin=*]
    \item \textbf{Number Theory}: Understanding that divisibility determines when horses return
    \item \textbf{Divisor Function}: For a given $n$, we need to count how many elements of $\{1,2,\ldots,10\}$ divide $n$
    \item \textbf{Prime Factorization}: $12 = 2^2 \cdot 3$ gives divisors $\{1, 2, 3, 4, 6, 12\}$ - five from our set
\end{itemize}

\subsection*{C. Cross-Domain Translation}
\begin{itemize}[leftmargin=*]
    \item \textbf{From Kinematics to Number Theory}: The problem initially appears to be about motion, but immediately translates to divisibility
    \item \textbf{From Periodic Functions to Divisors}: "Horse $k$ returns to start" $\Leftrightarrow$ "period $k$ divides time $t$"
\end{itemize}
\end{tacticalbox}

\begin{conversionbox}
\subsection*{A. Recognition Index}
\begin{itemize}[leftmargin=*]
    \item \textbf{Problem Signature}: "periodic motion", "all return to start", "least time"
    \item \textbf{Technique Match}: LCM for "all horses", divisibility for "at least some horses"
    \item \textbf{Recognition Time}: < 30 seconds - immediate identification of divisibility structure
\end{itemize}

\subsection*{B. High-Probability Techniques Applied}
\begin{itemize}[leftmargin=*]
    \item \textbf{Divisibility Analysis}: Primary technique - Horse $k$ at start $\Leftrightarrow$ $k \mid t$
    \item \textbf{Systematic Case Analysis}: Check $t = 1, 2, 3, \ldots$ until condition satisfied
    \item \textbf{Prime Factorization}: Helps identify which numbers have many divisors
\end{itemize}

\subsection*{C. Technique Integration Strategy}
\begin{enumerate}[leftmargin=*]
    \item \textbf{Step 1}: Recognize divisibility condition (1 min)
    \item \textbf{Step 2}: Systematically enumerate small $t$ values (1 min)
    \item \textbf{Step 3}: For each $t$, count divisors from $\{1,\ldots,10\}$ (1 min)
    \item \textbf{Step 4}: Find minimum $T$ with $\geq 5$ divisors (30 sec)
    \item \textbf{Step 5}: Compute digit sum (30 sec)
\end{enumerate}

\subsection*{D. Conflict Resolution}
\begin{itemize}[leftmargin=*]
    \item \textbf{Potential Conflict}: None - the divisibility approach is straightforward and unambiguous
    \item \textbf{Verification}: Can double-check by testing whether $T = 12$ indeed has at least 5 divisors from $\{1,\ldots,10\}$
\end{itemize}
\end{conversionbox}

\begin{verificationbox}
\subsection*{A. Verification Level Selection}
\begin{itemize}[leftmargin=*]
    \item \textbf{Selected Level}: Level 3 (Standard Verification with Alternative Method)
    \item \textbf{Justification}: Mid-band problem, 3 minutes available, need confidence but not exhaustive checking
\end{itemize}

\subsection*{B. Primary Verification Method}
\begin{itemize}[leftmargin=*]
    \item \textbf{Method}: Direct verification that $T = 12$ has exactly 5 divisors from $\{1,2,\ldots,10\}$
    \item \textbf{Execution}:
    \begin{itemize}
        \item Divisors of 12: $\{1, 2, 3, 4, 6, 12\}$
        \item From $\{1,\ldots,10\}$: $\{1, 2, 3, 4, 6\}$ - count = 5 $\checkmark$
        \item Check that no smaller number works:
        \begin{itemize}
            \item $t = 11$: divisors $\{1, 11\}$ - only 1 from our set
            \item $t = 10$: divisors $\{1, 2, 5, 10\}$ - only 4 from our set
            \item $t = 9$: divisors $\{1, 3, 9\}$ - only 3 from our set
            \item $t = 8$: divisors $\{1, 2, 4, 8\}$ - only 4 from our set
        \end{itemize}
    \end{itemize}
\end{itemize}

\subsection*{C. Alternative Method}
\begin{itemize}[leftmargin=*]
    \item \textbf{Method}: Use divisor formula - for $n = p_1^{a_1} \cdot p_2^{a_2} \cdots p_k^{a_k}$, total divisors = $(a_1+1)(a_2+1)\cdots(a_k+1)$
    \item \textbf{For $T = 12 = 2^2 \cdot 3$}: Total divisors = $(2+1)(1+1) = 6$
    \item \textbf{Count from $\{1,\ldots,10\}$}: All 6 divisors except 12 are $\leq 10$, so we have 5 divisors in our range $\checkmark$
\end{itemize}

\subsection*{D. Domain-Specific Checks}
\begin{itemize}[leftmargin=*]
    \item \textbf{Number Theory Check}: Verify $12$ is indeed the smallest number with at least 5 divisors $\leq 10$
    \item \textbf{Digit Sum Check}: $1 + 2 = 3$ matches answer choice (B) $\checkmark$
\end{itemize}

\subsection*{E. Confidence Calibration}
\begin{itemize}[leftmargin=*]
    \item \textbf{Initial Confidence}: 90\% (clear divisibility structure)
    \item \textbf{Post-Verification Confidence}: 98\% (verified with multiple methods)
    \item \textbf{Flag for Review}: No - high confidence after verification
\end{itemize}
\end{verificationbox}

\begin{roadmapbox}
\subsection*{A. Step-by-Step Solution Plan}

\textbf{Step 1: Understand the Divisibility Condition (30 seconds)}
\begin{itemize}[leftmargin=*]
    \item Recognize that Horse $k$ is at the starting point at time $t$ if and only if $k$ divides $t$
    \item Problem reduces to: Find minimum $T > 0$ such that at least 5 numbers from $\{1,2,\ldots,10\}$ divide $T$
\end{itemize}

\textbf{Step 2: Systematic Enumeration (1.5 minutes)}
\begin{itemize}[leftmargin=*]
    \item Test small values systematically:
    \begin{align*}
        t = 1: &\quad \text{divisors from } \{1,\ldots,10\} = \{1\} \quad (1 \text{ divisor}) \\
        t = 2: &\quad \{1, 2\} \quad (2 \text{ divisors}) \\
        t = 3: &\quad \{1, 3\} \quad (2 \text{ divisors}) \\
        t = 4: &\quad \{1, 2, 4\} \quad (3 \text{ divisors}) \\
        t = 5: &\quad \{1, 5\} \quad (2 \text{ divisors}) \\
        t = 6: &\quad \{1, 2, 3, 6\} \quad (4 \text{ divisors}) \\
        t = 7: &\quad \{1, 7\} \quad (2 \text{ divisors}) \\
        t = 8: &\quad \{1, 2, 4, 8\} \quad (4 \text{ divisors}) \\
        t = 9: &\quad \{1, 3, 9\} \quad (3 \text{ divisors}) \\
        t = 10: &\quad \{1, 2, 5, 10\} \quad (4 \text{ divisors}) \\
        t = 11: &\quad \{1\} \quad (1 \text{ divisor}) \\
        t = 12: &\quad \{1, 2, 3, 4, 6\} \quad (5 \text{ divisors}) \quad \checkmark
    \end{align*}
\end{itemize}

\textbf{Step 3: Verify Minimality (30 seconds)}
\begin{itemize}[leftmargin=*]
    \item Confirm no number less than 12 has 5 or more divisors from $\{1,\ldots,10\}$ (from enumeration above)
    \item Therefore, $T = 12$
\end{itemize}

\textbf{Step 4: Compute Digit Sum (15 seconds)}
\begin{itemize}[leftmargin=*]
    \item $T = 12$
    \item Digit sum = $1 + 2 = 3$
\end{itemize}

\textbf{Step 5: Select Answer (15 seconds)}
\begin{itemize}[leftmargin=*]
    \item Answer is $\boxed{\textbf{(B) } 3}$
\end{itemize}

\subsection*{B. Checkpoint Strategy}
\begin{itemize}[leftmargin=*]
    \item \textbf{Checkpoint 1} (after Step 1): Confirmed divisibility interpretation
    \item \textbf{Checkpoint 2} (after Step 2): Found $T = 12$ with 5 divisors
    \item \textbf{Checkpoint 3} (after Step 3): Verified minimality
    \item \textbf{Checkpoint 4} (after Step 4): Computed digit sum correctly
\end{itemize}

\subsection*{C. Alternative Approaches}
\begin{itemize}[leftmargin=*]
    \item \textbf{Approach 1}: Use divisor formula to identify candidates
    \begin{itemize}
        \item Numbers with many divisors: $2^a$, $3^b$, $2^a \cdot 3^b$
        \item Check: $8 = 2^3$ (4 divisors), $9 = 3^2$ (3 divisors), $12 = 2^2 \cdot 3$ (5 divisors from our set)
    \end{itemize}
    \item \textbf{Approach 2}: Answer choice elimination
    \begin{itemize}
        \item Digit sum 3 suggests numbers like 12, 21, 30, 111, etc.
        \item Check smallest: $12$ works!
    \end{itemize}
\end{itemize}

\subsection*{D. Comprehensive Pitfall Guide}

\textbf{Conceptual Pitfalls:}
\begin{itemize}[leftmargin=*]
    \item \textbf{Misinterpreting "at least 5"}:
    \begin{itemize}
        \item ERROR: Looking for exactly 5 horses
        \item FIX: Need $\geq 5$ horses, so 5 or more divisors
    \end{itemize}
    \item \textbf{Forgetting divisor 1}:
    \begin{itemize}
        \item ERROR: Not counting Horse 1 (which is always at start)
        \item FIX: Remember 1 divides every positive integer
    \end{itemize}
    \item \textbf{Confusion with LCM}:
    \begin{itemize}
        \item ERROR: Thinking we need LCM of 5 horses
        \item FIX: We need a number divisible by at least 5 horses, which is different from LCM
    \end{itemize}
\end{itemize}

\textbf{Computational Pitfalls:}
\begin{itemize}[leftmargin=*]
    \item \textbf{Miscounting divisors}:
    \begin{itemize}
        \item ERROR: For $t=12$, listing divisors as $\{1,2,3,4,6,12\}$ and counting 6
        \item FIX: Only count divisors $\leq 10$, so $\{1,2,3,4,6\}$ gives 5
    \end{itemize}
    \item \textbf{Arithmetic error in digit sum}:
    \begin{itemize}
        \item ERROR: Computing $1 + 2 = 2$ instead of 3
        \item FIX: Double-check simple arithmetic
    \end{itemize}
\end{itemize}

\textbf{Strategic Pitfalls:}
\begin{itemize}[leftmargin=*]
    \item \textbf{Testing too slowly}:
    \begin{itemize}
        \item ERROR: Writing out full divisor sets for each $t$
        \item FIX: Quickly identify divisors mentally or skip numbers unlikely to have many divisors (primes, prime powers $> 5$)
    \end{itemize}
    \item \textbf{Not verifying minimality}:
    \begin{itemize}
        \item ERROR: Finding 12 works but not checking if smaller values also work
        \item FIX: Systematically check all $t < 12$ or use the enumeration table
    \end{itemize}
\end{itemize}
\end{roadmapbox}

\begin{competitionbox}
\subsection*{A. Time Management}
\begin{itemize}[leftmargin=*]
    \item \textbf{Allocated Time}: 3 minutes (Problem 12 of 25 on AMC12)
    \item \textbf{Actual Time Breakdown}:
    \begin{itemize}
        \item Problem reading and setup: 30 seconds
        \item Systematic enumeration: 1.5 minutes
        \item Verification and digit sum: 1 minute
        \item Total: ~3 minutes $\checkmark$
    \end{itemize}
    \item \textbf{Time-Saving Strategies}:
    \begin{itemize}
        \item Skip obvious non-candidates (primes $> 5$, large prime powers)
        \item Focus on numbers with form $2^a \cdot 3^b$ (many small divisors)
    \end{itemize}
\end{itemize}

\subsection*{B. Problem Prioritization}
\begin{itemize}[leftmargin=*]
    \item \textbf{Accessibility}: High - clear structure, systematic approach
    \item \textbf{Domain Comfort}: Medium-High for students familiar with divisibility
    \item \textbf{Risk Assessment}: Low risk - straightforward enumeration, minimal computation
    \item \textbf{Strategic Decision}: Solve immediately if comfortable with number theory; otherwise skip and return
\end{itemize}

\subsection*{C. Risk Management}
\begin{itemize}[leftmargin=*]
    \item \textbf{Confidence Threshold}: 98\% (verified with multiple methods)
    \item \textbf{Flag for Review}: No - high confidence
    \item \textbf{Guess Strategy}: Not needed (high confidence in answer)
\end{itemize}

\subsection*{D. Abandonment Criteria}
\begin{itemize}[leftmargin=*]
    \item \textbf{Soft Abandon Trigger}: None - problem solved efficiently
    \item \textbf{Hard Abandon Trigger}: Would apply if unable to identify divisibility condition within 1 minute
    \item \textbf{Actual Status}: Completed successfully within time budget
\end{itemize}
\end{competitionbox}

\begin{keyinsightbox}
\textbf{The Core Insight}: Horse $k$ returns to the starting point at time $t$ if and only if $k$ divides $t$. This immediately transforms a kinematics problem into a number theory problem about counting divisors.

The minimum time for at least 5 horses is the smallest positive integer with at least 5 divisors from $\{1, 2, 3, \ldots, 10\}$. Systematic enumeration quickly reveals $T = 12$, which has exactly 5 such divisors: $\{1, 2, 3, 4, 6\}$.
\end{keyinsightbox}

\begin{warningbox}
\textbf{Don't Confuse Divisor Counting!}

When counting divisors of 12, we get $\{1, 2, 3, 4, 6, 12\}$ - that's 6 total divisors. However, we only count divisors $\leq 10$ (corresponding to our 10 horses), so we count $\{1, 2, 3, 4, 6\}$ = 5 divisors.

Always remember: we're counting horses (numbers 1-10), not all mathematical divisors!
\end{warningbox}

\begin{tipbox}
\textbf{Pattern Recognition for Speed}

Numbers with many small divisors typically have the form $2^a \cdot 3^b$ because 2 and 3 are the smallest primes. This is why 12 = $2^2 \cdot 3$ appears early in our search - it's the smallest number that's "highly composite" enough to have 5 divisors $\leq 10$.

In competition, once you see this pattern, you can jump directly to testing 6, 8, 9, 10, 12 rather than checking every single number!
\end{tipbox}

\section*{Final Answer}
The minimum time $T$ when at least 5 horses return to the starting point is $T = 12$ minutes.

The sum of the digits of $T$ is $1 + 2 = 3$.

\[\boxed{\textbf{(B) } 3}\]

\section*{Confidence Assessment}
\begin{itemize}[leftmargin=*]
    \item \textbf{Confidence Level}: 98\% - Verified with systematic enumeration and alternative divisor-counting method
    \item \textbf{Verification Applied}: Level 3 (Standard with Alternative) - confirmed all numbers $< 12$ have fewer than 5 divisors from $\{1,\ldots,10\}$
    \item \textbf{Flag for Review}: No - extremely high confidence, multiple verification methods agree
    \item \textbf{Time Spent}: ~3 minutes (on target for Problem 12)
\end{itemize}

\section*{Key Takeaways}
\begin{enumerate}
    \item \textbf{Translation is Key}: Recognize when a problem about periodic motion is really asking about divisibility
    \item \textbf{Systematic Approach}: For "find minimum" problems with small search space, enumeration is often fastest
    \item \textbf{Divisor Patterns}: Numbers of form $2^a \cdot 3^b$ tend to have many small divisors - prioritize testing these
    \item \textbf{Verification Matters}: Always check that no smaller value works when claiming a minimum
\end{enumerate}

\end{document}